\section{Results}
\label{sec:results}
\begin{table*}[!t]
\caption{Four-round graph endpoints, with five paired blocks per model. Errors are mean fractions. $Q$ counts correctness-audit queries per arm/block, excluding matching and evaluation labels; steps are cumulative updates (mean and range). $\Delta$ is audited-minus-unaudited error with a paired 95\% block-$t$ interval (negative favors auditing). Qwen3-4B contrasts pair corresponding blocks across runs with slightly different base predictions. Models are not pooled. Qwen3-1.7B is in Table~\ref{tab:qwen17-policy-endpoints}; one-step outcomes are separate (Table~\ref{tab:one-step-heldout}).}
\label{tab:audit-results}
\centering\small\setlength{\tabcolsep}{2.5pt}
\begin{tabular}{@{}llcccccc@{}}
\toprule
Policy & Arm & $Q$ & Steps & ID error & OOD error & \shortstack{$\Delta$ID\\(95\% CI)} & \shortstack{$\Delta$OOD\\(95\% CI)}\\
\midrule
\multicolumn{8}{@{}l}{\textbf{Llama-3.2-3B}}\\
Base & --- & 0 & 0 & 0.752 & 0.791 & --- & ---\\
Unaudited & R & 0 & 64 & 0.442 & 0.570 & --- & ---\\
Unaudited & S & 0 & 64 & 0.440 & 0.553 & --- & ---\\
Adaptive, B=16 & R & 64 & 54.2 (50--57) & 0.414 & 0.551 & $-0.028\;[-0.155,0.098]$ & $-0.019\;[-0.114,0.076]$\\
Adaptive, B=16 & S & 64 & 55.0 (44--63) & 0.380 & 0.524 & $-0.060\;[-0.144,0.024]$ & $-0.029\;[-0.107,0.050]$\\
\addlinespace
\multicolumn{8}{@{}l}{\textbf{Qwen3-4B}}\\
Base, unaudited & --- & 0 & 0 & 0.388 & 0.603 & --- & ---\\
Base, audited & --- & 0 & 0 & 0.390 & 0.604 & --- & ---\\
Unaudited & R & 0 & 64 & 0.269 & 0.501 & --- & ---\\
Unaudited & S & 0 & 64 & 0.266 & 0.502 & --- & ---\\
Adaptive, B=16 & R & 64 & 54.4 (49--61) & 0.282 & 0.525 & $0.013\;[-0.058,0.084]$ & $0.024\;[-0.045,0.092]$\\
Adaptive, B=16 & S & 64 & 55.0 (52--57) & 0.273 & 0.515 & $0.007\;[-0.108,0.121]$ & $0.013\;[-0.074,0.100]$\\
\bottomrule
\end{tabular}
\end{table*}

\begin{table*}[!t]
\caption{One-step graph held-out results from each run's own base/null. Pass@1 and paired $S-R$ entries report means $\pm$ descriptive 95\% block-$t$ half-widths. Target incidence is a fraction of all answers. ID contrasts are primary; OOD and target-rate contrasts are secondary, without multiplicity adjustment.}
\label{tab:one-step-heldout}
\centering\small\setlength{\tabcolsep}{3pt}
\begin{tabular}{@{}llllllll@{}}
\toprule
Model & Split & Null & Pass@1 $R$ & Pass@1 $S$ & Pass@1 $S-R$ & Target $R/S$ & Target $S-R$\\
\midrule
Qwen3-1.7B & ID & 0.4980 & $0.5028\pm0.0025$ & $0.5030\pm0.0017$ & $0.0002\pm0.0009$ & 0.1328/0.1327 & $-0.0001\pm0.0025$\\
Qwen3-1.7B & OOD & 0.3530 & $0.3530\pm0.0034$ & $0.3528\pm0.0038$ & $-0.0002\pm0.0058$ & 0.0602/0.0608 & $0.0006\pm0.0036$\\
Qwen3-4B & ID & 0.6145 & $0.6110\pm0.0024$ & $0.6104\pm0.0016$ & $-0.0006\pm0.0036$ & 0.2184/0.2193 & $0.0009\pm0.0042$\\
Qwen3-4B & OOD & 0.3920 & $0.3960\pm0.0025$ & $0.3982\pm0.0020$ & $0.0022\pm0.0006$ & 0.2614/0.2624 & $0.0010\pm0.0064$\\
Llama-3.2-3B & ID & 0.2455 & $0.2498\pm0.0019$ & $0.2507\pm0.0022$ & $0.0009\pm0.0028$ & 0.1293/0.1289 & $-0.0004\pm0.0029$\\
Llama-3.2-3B & OOD & 0.2070 & $0.2152\pm0.0028$ & $0.2162\pm0.0010$ & $0.0010\pm0.0029$ & 0.0722/0.0724 & $0.0002\pm0.0020$\\
\bottomrule
\end{tabular}
\end{table*}


Our empirical analysis comprises four-round graph self-training and
separate isolated one-step comparisons, each with five complete paired
blocks.

\paragraph{Matched starting selections.}
Table~\ref{tab:matching} certifies identical round-one prompt IDs, correct
responses, $48/16$ correct/error counts, and supervised token lengths.
The Qwen3-1.7B certificate is in Table~\ref{tab:qwen-matching};
Qwen3-4B controls are discussed in Appendix~\ref{app:qwen4b-four-round}.
The rate denominators exclude cut answers. Uniform error selection in $R$
still yields $9$--$13$ (Llama), $6$--$12$ (Qwen3-1.7B), or
$9$--$10$ (Qwen3-4B) \texttt{nonshortest}
responses among its 16 errors, versus 16 in $S$.
Auditing starts from these same selections; later checkpoint-dependent
pools are not jointly rate-matched.

\paragraph{Llama: both arms improve.}
Figure~\ref{fig:graph-four-round}(a--b) shows accuracy increasing from $0.248$ ID and
$0.209$ OOD to unaudited $R/S$ endpoints $0.558/0.560$ and $0.430/0.447$.
Paired $S-R$ accuracy differences are $0.002\pm0.105$ ID and
$0.017\pm0.081$ OOD (95\% block-$t$ intervals).
Both initial-round projections are negative ($-0.233/-0.260$), in contrast
to the opposite signs in the analytical construction. They describe
16-step round displacements rather than isolated one-step updates.

\paragraph{Finite-budget auditing in Llama.}
\label{sec:audit-results}
Audited $R/S$ reach $0.586/0.620$ ID and $0.449/0.476$ OOD accuracy.
The same-arm audited-minus-unaudited error contrasts in
Table~\ref{tab:audit-results} favor auditing in mean.
Each audited arm spends 64 queries; weighting changes retained volume,
ESS, and actual steps (Table~\ref{tab:audit-composition}).
Consequently, this comparison combines changes in data reliability,
training volume, and optimizer effort; it does not establish label-cost
savings. 

\paragraph{Composition, feasibility, and scope.}
Unaudited final target-error fractions are $0.147/0.183$ ID and
$0.090/0.118$ OOD for $R/S$, despite similar overall accuracy.
Auditing reduces mean ID target incidence, but OOD $S$ remains $0.120$:
suppression is not uniform. Appendix~\ref{app:h100-results} gives
all rounds, block endpoints, update diagnostics, original plots, and
cross-task feasibility checks. Intervals are descriptive and unadjusted,
not preregistered confirmatory claims. Arithmetic lacks exact-length
feasibility; complete costs remain unavailable.
Population illustrations are separate from these empirical findings.

\paragraph{Qwen3-4B: improvement without endpoint separation.}
Figure~\ref{fig:graph-four-round}(c--d) shows unaudited $R/S$ accuracy reaching
$0.7309/0.7339$ ID and $0.4988/0.4982$ OOD. Audited endpoints are
$0.7180/0.7273$ and $0.4750/0.4850$. Both policies improve in mean.
Same-arm audit error changes are positive
in mean (Table~\ref{tab:audit-results}).
Audited $R/S$ average 54.4/55.0 updates versus 64 unaudited;
these cross-run contrasts do not establish a corrective audit effect.
Qwen3-1.7B accuracy, composition, and audit results are retained in
Appendix~\ref{app:qwen-four-round}. Its audit accuracy changes are negative
in mean; the budget does not fix
optimizer effort. Neither model establishes a uniform audit benefit.
\paragraph{One-step held-out outcomes.}
Table~\ref{tab:one-step-heldout} reports each run against its own null;
run-specific baselines preclude a controlled cross-model ranking.
Qwen3-4B's secondary OOD contrast is $0.0022\pm0.0006$ (2--3 questions
per block), without multiplicity adjustment
(Figure~\ref{fig:one-step-contrasts}, Table~\ref{tab:one-step-heldout}).

\paragraph{Gradient--accuracy mismatch.}
With matched counts and shared correct responses, $S$ has larger mean
error-gradient contribution norms across models
(Table~\ref{tab:one-step-diagnostics}). Qwen3-4B has negative mean
$h^T\Delta\theta$ yet lower mean ID accuracy than null in both arms.
These diagnostics alone establish neither harmful updates nor reliable
accuracy separation.
